%% =============================================================================
%% SKELETON - Section 5: Factor Selection and Spanning Regressions
%% =============================================================================
%% Tables/figures INLINE (working draft for supervisor).
%% Same structure as paper1_sec5_factor_selection.tex.
%% =============================================================================


% ─────────────────────────────────────────────────────────────────────────────
\subsection{Principal Component Analysis}
\label{sec:pca}
% ─────────────────────────────────────────────────────────────────────────────

\subsubsection{Rolling PCA Methodology}
\label{sec:pca_method}

We extract principal components from a rolling window of $W=24$ months
over 80+ standardised candidate factors (\citealp{ludvigson2009macro}).
$K=8$ components are retained (approximately 80\% cumulative variance;
Figure~\ref{fig:pca_scree}). Both predictive ($PC_t \to R_{t+1}$)
and contemporaneous ($PC_t \to R_t$) timing are tested. Spanning
regressions are estimated by OLS with Newey--West HAC standard errors.
Robustness across $(W,K)$ grids is reported in
Appendix~\ref{app:pca_robustness}.

\begin{figure}[H]
\centering
\includegraphics[width=0.85\textwidth]{\figpath/pca/pca_scree_plot.pdf}
\caption{Scree Plot: Cumulative Variance Explained by Principal Components.
Dashed line indicates 80\% threshold. Baseline specification uses $K=8$ components.}
\label{fig:pca_scree}
\end{figure}


\subsubsection{Results}
\label{sec:pca_results}

Table~\ref{tab:pca_spanning} reports spanning regression results.
Alpha is statistically significant for all three strategies under both
timing conventions, while $\bar{R}^2$ remains moderate, indicating that
diffuse systematic risk does not span strategy returns.

\input{\tablepath/PCA_article}


% ─────────────────────────────────────────────────────────────────────────────
\subsection{Adaptive Elastic Net}
\label{sec:aen}
% ─────────────────────────────────────────────────────────────────────────────

\subsubsection{Candidate Factor Universe and Preprocessing}
\label{sec:aen_preproc}

The same 80+ candidate factors used in the PCA analysis
(Appendix~\ref{app:factor_list}) are preprocessed for penalised
regression: stationarity is verified via ADF tests, factors are
standardised to unit $L_2$-norm \citep{zou2009adaptive}, and
near-duplicate pairs ($|\rho| > 0.95$) are removed. After cleaning,
approximately 75+ factors remain per strategy.


\subsubsection{Estimation and Tuning}
\label{sec:aen_estimation}

The Adaptive Elastic Net \citep{zou2009adaptive} proceeds in two stages.
Stage~1 estimates a standard Elastic Net:
\[
\hat{\beta}^{EN} = \arg\min_{\beta}
\left\{
\frac{1}{T}\sum_{t=1}^{T}(r_t - \alpha - \beta'f_t)^2
+ \lambda_2\|\beta\|_2^2 + \lambda_1\|\beta\|_1
\right\}.
\]
Stage~2 constructs adaptive weights
$w_j = |\hat{\beta}^{EN}_j|^{-\gamma}$ ($\gamma=1$) and re-solves
with factor-specific $\ell_1$ penalties:
\[
\hat{\beta}^{AEN} = \arg\min_{\beta}
\left\{
\frac{1}{T}\sum_{t=1}^{T}(r_t - \alpha - \beta'f_t)^2
+ \lambda_2\|\beta\|_2^2
+ \lambda_1\sum_{j=1}^{p} w_j|\beta_j|
\right\}.
\]
Both stages use the same $\lambda_2$; $(\lambda_1, \lambda_2)$ are
selected by minimising the Hannan--Quinn criterion (HQC). The solution
is obtained via coordinate descent \citep{friedman2010regularization}.

\subsubsection{Stability Selection}
\label{sec:aen_stability}

To address selection uncertainty, we apply stability selection
\citep{meinshausen2010stability} adapted for time-series dependence:
$B=100$ stationary bootstrap resamples (expected block length 9 months)
are drawn, and the full AEN is re-run on each. A factor is retained
only if its selection frequency $\hat{\pi}_j \geq 60\%$. The expected
number of false selections is bounded by
$q^2 / [p(2\pi^* - 1)]$. Cluster-level frequency aggregation corrects
for vote-splitting among correlated factors
%(Figure~\ref{fig:stability_barplot}).

%\begin{figure}[H]
%\centering
%\includegraphics[width=\textwidth]{\figpath/aen/aen_stability_frequencies.pdf}
%\caption{Bootstrap Selection Frequencies for Candidate Risk Factors.
%Dashed line at $\hat\Pi = 60\%$ threshold. Blue bars indicate stably selected factors.}
%\label{fig:stability_barplot}
%\end{figure}


\subsubsection{Post-Selection OLS Results}
\label{sec:aen_results}

Once the stable set $\hat{S}_i$ is identified, we estimate unrestricted
OLS on the original (un-standardised) factors:
\[
r_{i,t} = \alpha_i + \beta_i' f_{\hat{S}_i,t} + \varepsilon_{i,t}
\]
with Newey--West HAC standard errors (6 lags).

Table~\ref{tab:aen_stable_ols} reports results. Alpha is statistically
significant for all three strategies. The selected factors are
economically interpretable and span several categories: liquidity and
funding frictions, credit risk premia, equity style factors, and
interest rate term structure variables. Each strategy loads on a
distinct combination, consistent with the different market
microstructures and investor bases involved. Compared with the three
benchmark frameworks (Section~4) and the PCA spanning regressions,
the AEN-selected factors deliver higher adjusted $\bar{R}^2$ while
preserving significant alpha, confirming that these targeted risk
drivers capture strategy-specific exposures that diffuse models miss.

\input{\tablepath/AEN_Stable_OLS_article}


% ─────────────────────────────────────────────────────────────────────────────
\subsection{Conditional Alpha}
\label{sec:conditional_alpha}
% ─────────────────────────────────────────────────────────────────────────────

Alpha increases during periods of financial stress across all strategies
and both factor-selection methods (Table~\ref{tab:alpha_regime_pca_aen}).
This is consistent with the limits-to-arbitrage mechanism
\citep{duffie2010presidential, mitchell2007slow}: when funding constraints
bind, capital retreats, mispricings widen, and alpha grows. Detailed
Ferson--Schadt and threshold robustness results are in
Appendix~\ref{app:pca_conditional} and~\ref{app:aen_conditional}.

\input{\tablepath/alpha_regime_synthesis_pca_aen}

% Ferson-Schadt and threshold robustness (AEN) in Appendix~\ref{app:aen_conditional}.
% Ferson-Schadt and threshold robustness (PCA) in Appendix~\ref{app:pca_conditional}.


% ─────────────────────────────────────────────────────────────────────────────
\subsection{Discussion, Method Comparison, and Robustness}
\label{sec:factor_discussion}
% ─────────────────────────────────────────────────────────────────────────────

Table~\ref{tab:aen_method_comparison} compares alpha across estimation
methods: AEN, Elastic Net, LASSO, Adaptive LASSO, and ALASSO-Chen. Alpha is robust to
the choice of method. AEN is preferred for its oracle property, its
ability to handle collinearity, and the stability selection step that
reduces post-selection bias. Sensitivity to the correlation threshold
$\rho$ and the adaptive exponent $\gamma$ is reported in
Appendix~\ref{app:aen_sensitivity}.

\input{\tablepath/AEN_Method_Comparison_article}